\section{Where the drivers fail, and why}

\begin{plainwords}
The proposed driver fails 15 of 200 drives: 12 collisions, 2 times leaving the road, 1 timeout. The Collision-Aware Teacher succeeds in \textbf{14 of those 15}, so the failures are not unsolvable scenes: they are things the student cannot copy. Eight of the 15 end within the first 12 seconds, and in seven of them every other model also collides; only the collision-aware teacher survives.
\end{plainwords}
\explainedfigure{fig_birds_eye_view_fidelity_before_collision}{0.97}{Gap to the other vehicle as drawn in the model's input, in the last steps before two collisions}{For two crashes of a bird's-eye-view model, the true distance between the car and the nearest vehicle (black) is compared with the distance that the 64 by 64 grid shows (red). A hollow circle marks steps at which the vehicle is not in the grid at all. These numbers are copied from the diagnosis note of 10 November (two replayed episodes).}{The oncoming vehicle in the left panel first appears in the grid only 0.9 seconds before impact. In the right panel the grid shows the crossing vehicle 1.7 metres farther away than it really is, and the vehicle disappears from the grid 0.1 seconds before the collision.}{The model's view of the other vehicles is late and coarse: vehicles are drawn as blocks without heading, and the grid ends 17.5 metres ahead. This is a limit of the representation that every model using this grid shares, and it is the leading explanation for the remaining failures.}{fig:fidelity}

\begin{itemize}
\item \textbf{The driver's view is late and coarse.} Vehicles are drawn as blocks without heading and the grid ends 17.5\,m ahead, while the teacher sees every vehicle up to 45\,m with exact velocities. This is an \emph{information gap}; it is in the teammates' own grid code on GitHub (\pth{bev_state/assembler.py}), so every model using the grid, including the two baselines, shares it.
\item \textbf{Dense traffic is unsolved}, even by the rule-based teachers (success 40--45\%), so no driver trained from them can do better; the main results are restricted to sparse traffic.
\item \textbf{The baselines fail differently:} the traffic-trained baseline leaves the road in 13.5\% of the drives (the proposed driver: 1.0\%), the signature of drifting after a small error with only five recordings to learn from.
\item \textbf{The map overlay has a price:} a constant error of 1\,m in the car's estimated position lowers the proposed driver's success from 0.925 to 0.70, and 2\,m to 0.275; a second, independent lane measurement repairs it (0.875--0.95), so redundancy, not noise training, is what protects it.
\end{itemize}

\section{Videos}
\begin{plainwords}
One scenario (number 5169) was recorded in four weather conditions with the front camera only, in smooth real time, for the two baselines, the proposed driver and the image-only driver. Both baselines collide in all four conditions; the proposed driver succeeds in all four. The image-only driver fails in all four (see the table in the technical report). Files are in \pth{report/videos/}.
\end{plainwords}
\explainedfigure{fig_seed5169_front_camera_frames}{0.97}{Frames of the front-camera videos of scenario 5169}{Each row is one model in one weather condition (clear, or rain of intensity 1.0); the columns are moments of the episode (steps 100, 200, 240 and the last frame; 10 steps = 1 second). The picture is the front-left camera of the teammates' rig with the weather augmentation applied. Videos for fog of intensity 1.0 and rain of intensity 0.6 are in the same folder.}{Both baselines collide with another vehicle at about step 250 (25 seconds). The proposed model does not collide and drives on until it reaches the destination at about step 700 (70 seconds). The Image-Only Driving Model collides at step 230 in clear weather and leaves the road after 88 steps in the heaviest rain. Heavy rain makes the picture much worse but does not change what the grid-based models do.}{This scenario was chosen \emph{because} the baselines fail on it, so it illustrates a typical failure but says nothing about how often it happens (the 200-episode results do). The grid-based models do not read these pixels; they read the simulator-built grid with weather applied as measured. Only the image-only model reads pixels, and it fails here in all four weather conditions.}{fig:frames}


\section{Comparison with the current GitHub repository}\label{sec:g-github}
\begin{plainwords}
GitHub (checked 12 November 2026; newest commit of 8 September 2026) contains the teammates' modules and exactly two model files, the two baselines. Everything produced in this project is on this computer only.
\end{plainwords}
\begin{center}\small
\begin{tabularx}{\textwidth}{@{}>{\raggedright\arraybackslash}p{3.0cm}>{\raggedright\arraybackslash}p{4.4cm}X@{}}\toprule
Component & On GitHub & In this project (local only)\\\midrule
Models & Two baseline checkpoints & 33 models, including the proposed driver and four groups of reinforcement-learning models\\
Training & Behaviour cloning from five recordings & Dataset aggregation, two teachers, caution, position-error and degraded-perception variants\\
Grid & Plain grid; map overlay ``unimplemented'' & Map overlay as a policy-side wrapper (the grid code is unmodified)\\
Weather & ``Preliminary tooling'', no measurement & Measured detector and depth degradation; 200-drive test\\
Reinforcement learning & ``Missing''; the traffic-trained baseline is recommended as the start & Four attempts, none from that checkpoint; all negative\\
Baseline success & 0.90 (9 of 10 drives) & 0.745 on 200 drives (interval 0.680 to 0.800)\\\bottomrule
\end{tabularx}
\end{center}
\begin{cautionbox}[Departures from rules written on GitHub]
\begin{itemize}[leftmargin=1.2em,itemsep=1pt]
\item The README says not to feed exact simulator vehicle positions to a learned model. The camera-free surrogate does (an ideal detector), degraded statistically where weather is studied.
\item The baselines were evaluated through that surrogate, not through the repository's evaluator.
\item Three teammate files carry local edits (an optional extra-input branch; defaults unchanged).
\item The first two reinforcement-learning attempts coupled rain to tyre grip, against the README's statement that weather does not change dynamics (later removed).
\end{itemize}
\end{cautionbox}

\section{What we can claim, and what we cannot}
\begin{center}\small
\begin{tabularx}{\textwidth}{@{}>{\raggedright\arraybackslash}p{5.4cm}X@{}}\toprule
We can claim (200 paired drives) & We cannot claim\\\midrule
The proposed driver beats both baselines and ties its teacher in sparse traffic. & That the map overlay (rather than the training method) causes the gain: untested at 200 drives.\\
Simulated fog and rain remove most detections but do not lower sparse-traffic success. & Any robustness or weather adaptation, or any result for physical rain or fog.\\
Four reinforcement-learning attempts failed, one significantly (mechanism hypothesised). & That reinforcement learning cannot work: only that these attempts did not.\\
Most remaining failures are imitation failures solved by a better-informed teacher. & Results for dense traffic, for the live camera pipeline, or between training seeds (one seed per model).\\\bottomrule
\end{tabularx}
\end{center}

\section{What to do next}\label{sec:g-next}
\begin{plainwords}
Seven possible papers exist. The best \emph{single} paper combines the first three: a thorough evaluation of training strategies (imitation against reinforcement learning), the measured weather model with its null result, and the information gap that limits the drivers.
\end{plainwords}
\begin{center}\small
\begin{tabularx}{\textwidth}{@{}>{\raggedright\arraybackslash}p{0.7cm}>{\raggedright\arraybackslash}p{4.3cm}>{\raggedright\arraybackslash}X>{\raggedright\arraybackslash}p{2.6cm}@{}}\toprule
\# & Direction & Evidence in hand, and what is missing & Cost\\\midrule
1 & Pitfalls of closed-loop evaluation & 0.975 on 40 drives became 0.895 on 200; validation of 1.0 hid a real loss; invalid first measurement; non-repeatable traffic. Missing: a resampling study of the selection bias. & Analysis of existing data\\
2 & Where does perception damage start to matter? & Measurement and the null result exist. Missing: a grid of detector loss against traffic density. & About 9 hours per model\\
3 & The teacher--student information gap & 14 of 15 failures solved by the collision-aware teacher. Missing: the unevaluated student of that teacher, a heading-aware grid, three training seeds. & 1 to 2 days\\
4 & Imitation against reinforcement learning & Six models at 200 drives, four documented failures. Missing: seeds, a stronger reinforcement-learning attempt. & About 2 hours per model\\
5 & Robustness to position error by redundancy & Curves at 40 drives. Missing: 200 drives, a real second cue. & About 11 hours\\
6 & Learned adaptive behaviour: caution from perception reliability & About 35\% less estimated energy at 200 drives, success tied at a 160-second limit. Missing: a regime where weather hurts, and a speed-matched control. & About 9 hours\\
7 & Benchmark and dataset release & Everything is stored with hashes. Missing: push to GitHub, licence check, repeatable traffic. & 1 to 2 days\\\bottomrule
\end{tabularx}
\end{center}
\textbf{Experiments to add first, in order:} (i) evaluate the student of the collision-aware teacher and the plain-grid six-number student on the 200 test drives (about 1.5 hours); (ii) two more training seeds for the proposed driver; (iii) the detector-loss against traffic-density grid; (iv) a heading-aware grid.

\section{Where to find everything}
\begin{center}\small
\begin{tabularx}{\textwidth}{@{}>{\raggedright\arraybackslash}p{5.2cm}X@{}}\toprule
What & Where (under \pth{D:/bev_Fresh/bevcapstone/})\\\midrule
This guide & \pth{report/Directors_Research_Guide.pdf}\\
Complete technical report (about 80 pages, every detail, file inventory) & \pth{report/Complete_Technical_Research_Report.pdf}\\
All 33 models, mapping of old to new names, hashes & \pth{trained_models/} (\pth{MODEL_NAME_MAPPING.md}, \pth{MODEL_MANIFEST.json})\\
Front-camera videos (scenario 5169) & \pth{report/videos/}\\
Figures and the scripts that produce them & \pth{report/figures/}, \pth{report/make\_report\_figures.py}\\
Code of this project (not on GitHub) & \pth{bev_rl/} (results in \pth{bev_rl/adaptive/results/})\\
Teammates' modules (on GitHub) & \pth{bev_policy/}, \pth{bev_state/}, \pth{camera/}, \pth{YOLO+weather module/}, \pth{baseline/}\\
Earlier drafts, first version of the report & \pth{archive/}\\\bottomrule
\end{tabularx}
\end{center}
